%%%%%%%%%%%%%%%%%%%%%%%
%%% DOCUMENT SETTINGS
%%%%%%%%%%%%%%%%%%%%%%%
\documentclass[11.0pt]{exam}
\printanswers % Uncomment to show answers
\include{formatAndDefs}
%%%%%%%%%%%%%%
%%% PACKAGES
%%%%%%%%%%%%%%
\usepackage{amsmath,amsfonts,amssymb,amsthm} % math fonts and symbols
\usepackage{bbm} % Math blackboard font
\usepackage{enumitem} % For reference enumerations lists
\usepackage{textcomp} % some symbols (like copyright)
\usepackage[top = 2cm, bottom = 2cm, right = 1.2cm, left = 1.2cm, paper = a4paper, headheight = 6cm]{geometry} % Margins setup
\usepackage{xcolor} % Colors
\usepackage{pgfplots}
\usepackage{spalign}
%%%%%%%%%%%%%%%%%%%%%%%%%%
%%% MACROS AND COMMANDS
%%%%%%%%%%%%%%%%%%%%%%%%%%
\definecolor{CUNEForange}{RGB}{237, 114, 1}
\newcommand{\N}{\mathbb{N}} % natural number set
\newcommand{\R}{\mathbb{R}} % real number set
%%%%%%%%%%%%%%%%%%%%%%%%%%%%%%%%%
%%% HEADER AND FOOT 
%%%%%%%%%%%%%%%%%%%%%%%%%%%%%%%%%
% Defining information
\newcommand{\degree}[1]{Bachelor in Philosophy, Politics and Economics}
\newcommand{\shortdeg}[1]{PPE}
\newcommand{\exam}[1]{Final Exam}
\newcommand{\subject}[1]{Mathematics}
\newcommand{\theyear}[1]{2023/24}
%\newcommand{\ccauthor}[1]{A. Azze \textcopyright\the\year}
\newcommand{\thedate}[1]{12 11 2023}
\newcommand{\university}[1]{CUNEF University}
% Setting header
\usepackage{fancyhdr}
\pagestyle{fancy}
\renewcommand{\headrulewidth}{0.1mm}
\renewcommand{\footrulewidth}{0.1mm}
% \renewcommand{\sectionmark}[1] % To avoid sections' titles being displayed on the header
{\markright{\MakeUppercase{}}}
\setlength\headsep{0.5cm}
% Header and footer
\fancyhead[L]{\degree{} \\ \subject{} - \theyear{}}
\fancyhead[R]{\thedate{} \\ \ccauthor{}}
\fancyfoot[L]{\university{}}
\fancyfoot[C]{}
\fancyfoot[R]{\thepage}

%%%%%%%%%%%%%%%%%%%%%%%%%
%%% DOCUMENT CONTENT
%%%%%%%%%%%%%%%%%%%%%%%%%
\begin{document}
%--- Instructions ---%
\begin{center}

    {\Huge \textbf{\exam{}}}
    \vspace{0.7cm}

    \begin{flushright}
        \textbf{Time Limit:} 120 minutes
    \end{flushright}
    \vspace{-0.2cm}
    
    \fbox{\parbox{\textwidth}{\centering 
        \begin{enumerate}[leftmargin = 0.5cm, label = $\bullet$, rightmargin = 0.2cm]
    
            \item \textbf{Every non-trivial calculation and claim in the exam must be properly justified.}

            \item The use of calculators is not necessary and therefore not allowed.

            \item Digital devices such as mobile phones, tablets, and laptops, as well as internet access, are forbidden.

            \item You can use physical resources such as books, notes, and other printed materials to assist you during the exam.
 
        \end{enumerate}
    }}
\end{center}
%--- Instructions ---%

%--- Questions ---%

\begin{questions}

% Solution
%\begin{solution}[0cm]
%    \begin{enumerate}[leftmargin = 0.5cm, label = (\alph*), rightmargin = 0.2cm]
%\end{solution}

% Question
\question[3] Determine if the following functions have global minimum and/or maximum and compute them.
\begin{enumerate}
    \item $f(x)=2x^{3}+3x^2-12x$ in the interval $[-2,2]$.
    \item $f(x)=2x^{3}+3x^2-12x$ in the interval $(-2,2)$.
    \item $f(x,y)=x^2-2xy-y^2-12x+16y$ in $\mathbb{R}$.
\end{enumerate}

% Solution
\begin{solution}[0cm]
\begin{enumerate}
    \item The function is continuous and the interval is closed, so we can assure, by the Extreme Value Theorem, that there will be minimum and maximum. The derivative of $f(x)$ is $6x^2+6x-12$ which is zero when $x=-2$ and $x=1$. We evaluate the function on these points and the extremes of the interval (the critical points), obtaining $f(-2)=20, f(1)=-7, f(2)=4$, hence the minimum is $-7$ and reached at $x=1$ and the maximum is 20 and reached at $x=-2$.
    \item In this case, we have no guaranteed the existence of minimum and maximum since $(-2,2)$ is not closed. But, if the minimum in $[-2,2]$ was in $x=1$, the same holds for $(-2,2)$. For the maximum, as it was in the extreme $x=-2$ that now is excluded, but very close points to $-2$ will also have larger values than the rest of the points in the interval, but without a limit, so there is no maximum. We can study the growth of the function to get that it is as shown in the following figure. So, the maximum would be at $x=-2$ but we never reach it, since $x=-2$ is not included.
\begin{center}
 \begin{tikzpicture}
  \begin{axis}[
    axis lines=middle,
    xlabel=$x$,
    ylabel=$y$,
    xmin=-2.5,
    xmax=2.5,
    ymin=-20,
    ymax=20,
    xtick={-2, -1, 0, 1, 2},
    ytick={-20, -15, -10, -5, 0, 5, 10, 15, 20},
    samples=100,
    domain=-2:2,
    ]
    \addplot[blue, thick] {2*x^3 + 3*x^2 - 12*x};
  \end{axis}
\end{tikzpicture}
\end{center}
\item To get the critical points of $f(x,y)$ we compute the partial derivatives, obtaining \[
   f_x=2x-2y-12,\enspace f_y=-2x-2y+16.
\] The critical points, resulting from solving $f_x=f_y=0$ are just the point $(7,1)$. To determine if we have a minimum, maximum or a saddle point, we use the second derivative criterion. We compute the second order derivatives, getting \[f_{xx}=2,\enspace f_{yy}=-2,\enspace f_{xy}=-2,\] an hence $H(x,y)=2(-2)-(-2)^2=-8$ and hence the critical point is just a saddle point. The function has neither a minimum nor a maximum.
\end{enumerate}
\end{solution}

% Question
\question[3]
Compute the following integrals.
\begin{enumerate}
    \item $\displaystyle \int (1+x^2) \sin(x)\, dx$
    \item $\displaystyle \int \frac{x+1}{(x+3)(x-2)} \, dx$
    \item $\displaystyle \int_1^{\infty} xe^{-5x^2} \, dx$
\end{enumerate}

\begin{solution}[0cm]
\begin{enumerate}
\item To solve the integral

\[
I=\int (1 + x^2) \sin(x) \, dx
\]

We apply integration by parts. Let's define \( u \) and \( dv \) as

\[
\begin{aligned}
u &= (1 + x^2), \quad &dv &= \sin(x) \, dx \\
du &= 2x \, dx, \quad &v &= -\cos(x)
\end{aligned}
\]

Applying the integration by parts formula, we get

\[
\begin{aligned}
\int (1 + x^2) \sin(x) \, dx &= uv - \int v \, du =\\
&= -(1 + x^2)\cos(x) - \int -\cos(x) \cdot 2x \, dx=\\
&= -(1 + x^2)\cos(x) + 2 \int x \cos(x) \, dx \\
\end{aligned}
\]

We apply integration by parts again. Let's define \( u \) and \( dv \) as

\[
\begin{aligned}
u &= x, \quad &dv &= \cos(x) \, dx \\
du &= dx, \quad &v &= \sin(x)
\end{aligned}
\]
We get 

\[
\begin{aligned}
I &= -(1 + x^2)\cos(x) + 2\left(x\sin(x) - \int \sin(x) \, dx\right)=\\
&= -(1 + x^2)\cos(x) + 2x\sin(x) - 2\cos(x) + C=\\
&= - x^2\cos(x) + 2x\sin(x) - \cos(x) + C.
\end{aligned}
\]

\item To solve the integral \[
I=\int \frac{x + 1}{(x + 3)(x - 2)} \, dx
\]
we can use partial fraction decomposition. Write the function as

\[
\frac{x + 1}{(x + 3)(x - 2)}=\frac{A}{x + 3} + \frac{B}{x - 2}=\frac{A(x-2)+B(x+3)}{(x + 3)(x - 2)}=\frac{Ax-2A+Bx+3B}{(x + 3)(x - 2)}
\]
and hence we have
\[
\begin{aligned}
A+B&=1\\
-2A+3B&=1
\end{aligned}
\]
which gives us $A=\frac{2}{5}, B=-\frac{2}{5}$. Hence, the integral is
\[I=\int \frac{\frac{2}{5}}{x+3}+\frac{-\frac{3}{5}}{x-2}\, dx = \frac{2}{5}\ln(x+3)-\frac{2}{5}\ln(x-2)+C.\]

\item The integral \[
\int_1^{\infty} xe^{-5x^2} \, dx
\]
converges if so does the limit
\[\lim_{b\rightarrow\infty}\int_1^b xe^{-5x^2} \,dx.\]
We compute the integral 
\[I=\int_1^b xe^{-5x^2}=\frac{-e^{-5b^2}}{10}\, dx= \frac{-e^{-5x^2}}{10}+\frac{e^{-5}}{10}.\]
Finally, we compute the limit
\[\lim_{b\rightarrow\infty} I=\frac{e^{-5}}{10}.\]
Hence, the original integral converges to $\displaystyle \frac{1}{10e^5}$.
\end{enumerate}
\end{solution}
% Question
\question[4] Consider the matrices
\begin{equation*}
A=\begin{pmatrix}
-1 & 0 & 0\\
0 & 1 & 0\\
0 & 0 & -1
\end{pmatrix},
\enspace
B=\begin{pmatrix}
1 & 1 & 0\\
2 & -1 & 3\\
3 & 0 & 3
\end{pmatrix},
\enspace
t=\begin{pmatrix}
x\\
y\\
z
\end{pmatrix},
\enspace
c=\begin{pmatrix}
1\\
2\\
0
\end{pmatrix}.
\end{equation*}
\begin{enumerate}
\item Compute $A^{10}$.
\item Compute the determinant of $B$. What is the inverse of $B$?
\item Determine the value of the matrix $C$ that satisfies $(A+B)C=I$, where $I$ is the identity matrix of size 3.
\item Use the Gaussian elimination procedure to solve the linear system $Bt=c$
\end{enumerate}

\begin{solution}
\begin{enumerate}
    \item It is very easy to see that, the powers of a diagonal matrix consist of the powers of its elements, that is 
\[A=\begin{pmatrix}
(-1)^{10} & 0 & 0\\
0 & 1^{10} & 0\\
0 & 0 & (-1)^{10}
\end{pmatrix}=
\begin{pmatrix}
1 & 0 & 0\\
0 & 1 & 0\\
0 & 0 & 1
\end{pmatrix}.\]
\item The determinant of $B$ is $|B|=-3+9-6=0$, hence the matrix $B$ has no inverse $B^{-1}$.
\item To compute the matrix $C$, we have that 
\[A+B=\begin{pmatrix}
0 & 1 & 0\\
2 & 0 & 3\\
3 & 0 & 2
\end{pmatrix},\]
so we have to isolate $C$ from
\[
\begin{pmatrix}
0 & 1 & 0\\
2 & 0 & 3\\
3 & 0 & 2
\end{pmatrix}C=
\begin{pmatrix}
1 & 0 & 0\\
0 & 1 & 0\\
0 & 0 & 1
\end{pmatrix}\]
and we can do it left-multiplying by the inverse of $A+B$ in both sides,
\[
\begin{pmatrix}
0 & 1 & 0\\
2 & 0 & 3\\
3 & 0 & 2
\end{pmatrix}^{-1}\cdot
\begin{pmatrix}
0 & 1 & 0\\
2 & 0 & 3\\
3 & 0 & 2
\end{pmatrix}C=
\begin{pmatrix}
0 & 1 & 0\\
2 & 0 & 3\\
3 & 0 & 2
\end{pmatrix}^{-1}\cdot
\begin{pmatrix}
1 & 0 & 0\\
0 & 1 & 0\\
0 & 0 & 1
\end{pmatrix}\]
which gives us
\[C=\begin{pmatrix}
0 & 1 & 0\\
2 & 0 & 3\\
3 & 0 & 2
\end{pmatrix}^{-1}.\]
The determinant of this matrix is 5 and its cofactor matrix is 
\[\begin{pmatrix}
0 & 5 & 0\\
-2 & 0 & 3\\
3 & 0 & -2
\end{pmatrix}.\]
Hence, the inverse is
\[C=\frac{1}{5}\begin{pmatrix}
0 & -2 & 3\\
5 & 0 & 0\\
0 & 3 & -2
\end{pmatrix}.\]
\item The system $Bt=c$ is given by the equations 
\[
  \spalignsys{
x+y=1 ;
2x-y+3z=2 ;
3x+3z=0
}
\]
We apply the Gaussian elimination procedure to the matrix
\[
\left(\begin{array}{ccc|c}
    1 & 1 & 0 & 1 \\
    2 & -1 & 3 & 2 \\
    3 & 0 & 3 & 0 \\
\end{array}\right)
\]

The goal is to transform this matrix into the echelon form. The steps are ($R2$ means row 2):
\[
\begin{aligned}
R2&\rightarrow R2-2R1\\
R3&\rightarrow R3-3R1\\
R3&\rightarrow R3-R2
\end{aligned}
\]
This gives the following equivalences.
\[
\left(\begin{array}{ccc|c}
   1 & 1 & 0 & 1 \\
    2 & -1 & 3 & 2 \\
    3 & 0 & 3 & 0 \\
\end{array}\right)\sim
\left(\begin{array}{ccc|c}
   1 & 1 & 0 & 1 \\
    0 & -3 & 3 & 0 \\
    3 & 0 & 3 & 0 \\
\end{array}\right)\sim
\left(\begin{array}{ccc|c}
    1 & 1 & 0 & 1 \\
    0 & -3 & 3 & 0 \\
    0 & -3 & 3 & -3 \\
\end{array}\right)\sim
\left(\begin{array}{ccc|c}
    1 & 1 & 0 & 1 \\
    0 & -3 & 3 & 0 \\
    0 & 0 & 0 & -3 \\
\end{array}\right).
\]
Thus, the original system is equivalent to the system
\[
  \spalignsys{
x+y=1 ;
-3y+3z=0;
0=-3
}
\]
which is clearly inconsistent, so there is no solution.
\end{enumerate}
\end{solution}
% % Question
% \question[2]
% The tobacco market is a monopoly, with supply and demand sets
% \begin{align*}
%     S = \{(q, p) : 15p − 2q = 20\},\quad D = \{(q, p) : q + 5p = 40\}.
% \end{align*}
% Given that the government has imposed an excise tax of $T$ dollars to the tabaco market:

% \begin{enumerate}[leftmargin = 0.5cm, label = (\alph*), rightmargin = 0.2cm]
        
%     \item Define the governments revenue function in terms of the excise tax and the number of units sold.

%     \item Find the excise tax that maximizes the government revenue.
        
% \end{enumerate}
% % Solution
% \begin{solution}[0cm]
%     Solution
% \end{solution}

% % Question
% \question[2]
%     A firm manufactures two goods, $X$ and $Y$. The market price of these goods is unaffected by the level of the firm's production. The firm's cost function is $C(x, y) = 2x^2+ xy + 2y^2$. The market prices for $X$ and $Y$ are $12$ and $18$ dollars per unit, respectively. Determine the number of units of each good that the firm should produce to maximize its profit.  
% % Solution
% \begin{solution}[0cm]
%     Solution
% \end{solution}

% % Question
% \question[2]
%     A wine dealer has certain amount of units of barrels of wine, which he can sell today at $K$ dollars or wait for a better deal. The price of his barrels of wine increases in time according the the function $Kt^n$, for $n \in \N$. Assume that the market holds an interest rate of $r \R_+$ on continuously compounding levels, that is, $x$ today's dollars equal $xe^{-r t}$ dollars $t$ units of time ahead.
    
%     \begin{enumerate}[leftmargin = 0.5cm, label = (\alph*), rightmargin = 0.2cm]

%         \item Define the profit made by selling the barrels of wine at time $t$.
    
%         \item When should the dealer sell the barrels of wine to maximize his profit? Hint: you can use Problem 3 to ease the optimization.
    
%     \end{enumerate}
% % Solution
% \begin{solution}[0cm]
%     Solution
% \end{solution}

% Question


\end{questions}	

\end{document}